\section*{अध्याय \ref{ch:g12:human-evolution} --- मानव विकास पर एक नज़र}

\begin{solution}{exo:g12:human-evolution:1}
चिंपैंज़ी (बोनोबो समेत); और वे वंशरेखाएँ छह से सात मिलियन साल पहले अलग
हुईं।
\end{solution}

\begin{solution}{exo:g12:human-evolution:2}
खोपड़ी के नीचे मेरु छिद्र; दोहरे मोड़ वाली रीढ़; छोटी, कटोरे जैसी श्रोणि;
घुटने की ओर भीतर को झुकी जाँघ की हड्डियाँ; और सीध में पड़े अंगूठे वाला
मेहराबदार पैर।
\end{solution}

\begin{solution}{exo:g12:human-evolution:3}
लगभग \qty{420}{cm^3}; लगभग \qty{850}{cm^3}; लगभग \qty{1350}{cm^3}।
\end{solution}

\begin{solution}{exo:g12:human-evolution:4}
क्योंकि उस वंशरेखा की कई \omterm{def:g10:biodiversity-scales:species}{जातियाँ} एक ही समय जीं, और उनमें से अधिकांश अगली
में बदलने के बजाय विलुप्ति पर ख़त्म हुईं; वे \omterm{def:g10:biodiversity-scales:species}{जातियाँ} शाखाएँ हैं, सीढ़ियाँ
नहीं।
\end{solution}

\begin{solution}{exo:g12:human-evolution:5}
अफ़्रीका में, लगभग 300\,000 साल पहले। \omterm{def:g10:biodiversity-scales:biodiversity}{आनुवंशिक विविधता} अफ़्रीका में सबसे
अधिक है और उससे दूरी के साथ गिरती है: सारे जीवित मनुष्य उसी अफ़्रीकी \omterm{def:g12:selection-drift-speciation:population}{आबादी}
से, एक के बाद एक छोटे संस्थापक समूहों के रास्ते, उतरे हैं।
\end{solution}

\begin{solution}{exo:g12:human-evolution:6}
\omterm{def:g12:human-evolution:lineage}{मानव वंशरेखा} में दो पूर्वज \omterm{def:g11:cell-cycle-mitosis:chromosome}{गुणसूत्र} जुड़कर एक हो गए, इसलिए 24 जोड़े 23 बन
गए। वह जुड़ा \omterm{def:g11:cell-cycle-mitosis:chromosome}{गुणसूत्र} बीच में दोनों मूल सिरों के अवशेष और एक दूसरा, चुप
\omterm{def:g11:cell-cycle-mitosis:chromosome}{सेंट्रोमियर} ढोता है --- यानी उस जुड़ाव का एक अभिलेख; और उसकी \omterm{def:g10:universal-dna:gene}{जीन} सामग्री
दोनों चिंपैंज़ी \omterm{def:g10:universal-dna:information}{गुणसूत्रों} से \omterm{def:g10:universal-dna:gene}{जीन} दर \omterm{def:g10:universal-dna:gene}{जीन} मेल खाती है। \omterm{def:g10:universal-dna:information}{गुणसूत्रों} के किसी
साझा सेट से, जो एक बार बदला हो, उतरना ठीक वही है जिसकी नातेदारी भविष्यवाणी
करती है।
\end{solution}

\begin{solution}{exo:g12:human-evolution:7}
सीधा चलना, और वह भी लगभग दो मिलियन साल से: 3.6 से 3.2 मिलियन साल पुराने
पैरों के निशान तथा श्रोणियाँ \qty{400}{cm^3} मस्तिष्क वाले चलने वालों की
हैं; और मस्तिष्क की बढ़त लगभग 2.5 मिलियन साल पर शुरू होती है।
\end{solution}

\begin{solution}{exo:g12:human-evolution:8}
दो मिलियन साल पहले: \emph{A. africanus} (मुश्किल से), \emph{Paranthropus},
\emph{Homo habilis} और शुरुआती \emph{H. erectus}। 100\,000 साल पहले:
\emph{H. erectus} (आख़िरी \omterm{def:g12:selection-drift-speciation:population}{आबादियाँ}), \emph{H. heidelbergensis} (आख़िरी),
निएंडरथल और \emph{H. sapiens}।
\end{solution}

\begin{solution}{exo:g12:human-evolution:9}
अफ़्रीका से निकलते आधुनिक मनुष्यों और पश्चिमी एशिया तथा यूरोप में मिले
निएंडरथलों के बीच कोई 50\,000 साल पहले हुआ अंतःप्रजनन। जो \omterm{def:g12:selection-drift-speciation:population}{आबादियाँ}
उप-सहारा अफ़्रीका में रह गईं वे निएंडरथलों से कभी मिलीं ही नहीं, इसलिए वे
कुछ भी नहीं ढोतीं।
\end{solution}

\begin{solution}{exo:g12:human-evolution:10}
सीधी चलती (आगे का छिद्र), छोटा मस्तिष्क, बड़े जबड़े, 2.8 मिलियन साल: यानी
कोई ऑस्ट्रैलोपिथेसीन, जो उस वंशरेखा के आधार के पास किसी बग़ल की शाखा पर
रखा जाएगा --- यानी \emph{होमो} के पूर्वजों का कोई नातेदार, कोई सिद्ध
पूर्वज नहीं।
\end{solution}

\begin{solution}{exo:g12:human-evolution:11}
महावानर एक \omterm{prop:g12:phylogenetic-trees:rules}{क्लेड} बनाते हैं, और मनुष्य उसके भीतर की एक शाखा हैं; और चूँकि
कोई \omterm{prop:g12:phylogenetic-trees:rules}{क्लेड} अपनी गाँठ के सारे वंशज समेटता है, इसलिए मनुष्य वानर हैं, जैसे वे
नर-वानर और स्तनधारी हैं। ``वानरों से उतरे'' ग़लत ढंग से वानरों को किसी गाँठ
पर बिठा देता है; जबकि जीवित वानर सहोदर सिरे हैं।
\end{solution}

\begin{solution}{exo:g12:human-evolution:12}
परिकल्पना 1: प्रतिस्पर्धा --- आधुनिक मनुष्यों ने, अधिक संख्या में या बेहतर
संगठित होकर, संसाधन ले लिए; और साक्ष्य वहाँ निएंडरथल स्थलों की गिरावट होती
जहाँ दोनों एक-दूसरे पर पड़े। परिकल्पना 2: समावेश --- छोटी निएंडरथल \omterm{def:g12:selection-drift-speciation:population}{आबादियाँ}
आंशिक रूप से अंतःप्रजनन से समेट ली गईं और आंशिक रूप से बदलती जलवायु में
अपवाह तथा संयोग से विलुप्त हुईं; और साक्ष्य हममें बैठा निएंडरथल \omterm{def:g10:universal-dna:information}{DNA} तथा
संपर्क से पहले छोटी, घटती निएंडरथल \omterm{def:g12:selection-drift-speciation:population}{आबादियों} के चिह्न होते। दोनों सच हो सकती
हैं।
\end{solution}

\begin{solution}{exo:g12:human-evolution:13}
उस प्रकीर्णन का हर क़दम किसी छोटे समूह ने उठाया जो उस \omterm{def:g12:selection-drift-speciation:population}{आबादी} के
\omterm{def:g10:universal-dna:gene}{युग्मविकल्पियों} का एक प्रतिदर्श ढोता था जिसे वह छोड़ आया था; और हर नया
इलाक़ा पिछले से बसाया गया। हर बसावट पर विविधता गिरी, इसलिए सबसे दूर की
\omterm{def:g12:selection-drift-speciation:population}{आबादियों} के पास सबसे कम है। वह प्रकीर्णन किसी सामूहिक हलचल के बजाय एक के
बाद एक छोटे प्रवासों से आगे बढ़ा।
\end{solution}

\begin{solution}{exo:g12:human-evolution:14}
दुग्ध-पालन ने दूध को वयस्कों का भोजन बनाया: टिकाऊपन वाले \omterm{def:g10:universal-dna:gene}{युग्मविकल्पी} के
वाहक बेहतर खाते थे और अधिक वंशज छोड़ते थे। कमज़ोर धूप वाले अक्षांशों में
प्रवास ने विटामिन D दुर्लभ कर दिया: इसलिए हल्की त्वचा, जो उसे अधिक बनाती
है, पसंद की गई। खेती ने घनी \omterm{def:g12:selection-drift-speciation:population}{आबादियाँ} और ठहरा पानी बनाया जहाँ मलेरिया फैला:
इसलिए उससे बचाते \omterm{def:g10:universal-dna:gene}{युग्मविकल्पी} चढ़े। हर हाल में जीने के ढंग के किसी बदलाव ने
यह बदल दिया कि कौन-से \omterm{def:g10:universal-dna:gene}{युग्मविकल्पी} प्रजनन करते हैं।
\end{solution}

\begin{solution}{exo:g12:human-evolution:15}
वरण अलग-अलग \omterm{def:g10:universal-dna:gene}{युग्मविकल्पी} ढोते वाहकों के बीच प्रजनन का कोई भी अंतर है; और
दवा यह बदलती है कि कौन-से \omterm{def:g10:universal-dna:gene}{युग्मविकल्पी} मायने रखते हैं, यह नहीं कि कोई रखता
है या नहीं --- संक्रमणों के प्रति प्रतिरोध, उर्वरता, प्रजनन की उम्र आज भी
बदलते हैं और विरासत में मिलते हैं। वे तीनों लक्षण पिछले दस हज़ार सालों के
भीतर काम करता वरण दिखाते हैं, जो \omterm{prop:g12:diversification-of-life:behaviour}{संस्कृति} ने ही चलाया। और अपवाह हर \omterm{def:g12:selection-drift-speciation:population}{आबादी} में
चाहे जो हो जारी रहता है। मानव विकास ने दिशा बदली है, रुका नहीं है।
\end{solution}

\begin{solution}{pb:g12:human-evolution:1}
\textbf{1.} A: $21.5/0.15 \approx \qty{1.43}{m}$; B: $18.5/0.15
\approx \qty{1.23}{m}$।

\textbf{2.} $0.8 \times 1.43 = \qty{1.14}{m}$ और $0.8 \times 1.23 =
\qty{0.98}{m}$: दोनों डग किसी धीमी चाल की डग से छोटी हैं; वे चल रहे थे,
बिना जल्दबाज़ी के।

\textbf{3.} गहरी एड़ी: यानी एड़ी पहले पड़ती है, जैसा मानव चलने में; एक
मेहराब: यानी ऐसा पैर जो उत्तोलक और स्प्रिंग की तरह काम करता है; समांतर
अंगूठा: यानी धक्का देने वाला पैर, टहनियाँ पकड़ने वाला नहीं।

\textbf{4.} \qty{27}{m} तक टिकी, एड़ी-पहले वाली डग भरती चाल, और वह भी मनुष्य
जैसे पैर के साथ: यानी आदतन द्विपादता। निशानों से अनजाना: श्रोणि, जाँघ की
हड्डी का कोण, मेरु छिद्र, और रीढ़ के मोड़।

\textbf{5.} यह कि सीधा चलना मस्तिष्क के बड़े होने से एक मिलियन साल से भी
अधिक पहले आया।

\textbf{6.} 1: ऑस्ट्रैलोपिथेसीन। 2: \emph{Homo habilis}। 3: \emph{Homo
erectus}। 4: निएंडरथल। 5: \emph{Homo sapiens}।

\textbf{7.} 3.0 मिलियन पर 450 से 1.9 पर 650, 1.0 पर 950, फिर 1450 और
1350: सबसे तेज़ बढ़त 1.9 और 1.0 मिलियन साल के बीच है (\qty{300}{cm^3} प्रति
मिलियन साल), फिर 1.0 और 0.06 के बीच।

\textbf{8.} नहीं: निएंडरथल मस्तिष्क उतने ही बड़े या बड़े थे। आधुनिक मनुष्य
ऊँचे माथे, कपाल-कोष्ठ के नीचे दुबके छोटे चेहरे, ठुड्डी और गोल खोपड़ी से
पहचाने जाते हैं।

\textbf{9.} उस झाड़ी की दो \omterm{def:g12:selection-drift-speciation:population}{आबादियाँ} जो अग़ल-बग़ल जीं, शरीर-रचना में अलग
थीं, और जिन्होंने अंतःप्रजनन किया: अफ़्रीका के बाहर जीवित लोगों में बैठा
निएंडरथल \omterm{def:g10:universal-dna:information}{DNA} ही उसका साक्ष्य है।

\textbf{10.} कोई जीवाश्म अपनी ही बग़ल की शाखा पर एक सिरा है; उसके लक्षणों
का जोड़ उसे उस गाँठ के पास रखता है जिससे बाद की \omterm{def:g10:biodiversity-scales:species}{जातियाँ} उतरती हैं, पर कोई
लक्षण यह साबित नहीं करता कि वह ठीक वही \omterm{def:g12:selection-drift-speciation:population}{आबादी} थी जिसने उन्हें जन्म दिया। वह
उस पूर्वज के पास का एक नातेदार है, और किसी जीवाश्म को इतना ही दिखाया जा
सकता है।

\textbf{11.} उसके जीनोम का 1.6\% बनाते खंड किसी भी आधुनिक अनुक्रम से
ज़्यादा निएंडरथल अनुक्रम के पास हैं: वे उन निएंडरथल पूर्वजों से विरासत में
मिले जिन्होंने आधुनिक मनुष्यों से अंतःप्रजनन किया।

\textbf{12.} वह चार्ट यूरोप में लगभग 1.8\% निएंडरथल और मेलानेशिया में उस
एशियाई \omterm{def:g12:selection-drift-speciation:population}{आबादी} का लगभग 3.5\% \omterm{def:g10:universal-dna:information}{DNA} देता है: उसके यूरोपीय वंश ने निएंडरथल हिस्सा
लाया, और उसके मेलानेशियाई वंश ने दूसरा (जो उसके मिश्रित वंशक्रम से घुल गया
है)।

\textbf{13.} $50\,000/25 = 2000$ पीढ़ियाँ। वह हिस्सा इसलिए नहीं घुला
क्योंकि पूरी \omterm{def:g12:selection-drift-speciation:population}{आबादी} उसे ढोती है: दूसरे वाहकों से मिलना उस औसत को स्थिर रखता
है; और घुलाव तभी होता है जब वाहक ग़ैर-वाहकों से मिलें, तथा अफ़्रीका के बाहर
कोई ग़ैर-वाहक थे ही नहीं।

\textbf{14.} यह कि वह सीमा बंद नहीं थी: कोई 500\,000 सालों के अलगाव के बाद
भी वे \omterm{def:g12:selection-drift-speciation:population}{आबादियाँ} उर्वर संतान दे सकती थीं --- यानी बनती हुई \omterm{def:g10:biodiversity-scales:species}{जातियाँ}, जो जीव
विज्ञान से ज़्यादा भूगोल से अलग-थलग थीं।

\textbf{15.} \omterm{def:g10:cells-common-unit:organelle}{सूत्रकणिकाएँ} केवल माँओं से जाती हैं: जीवित मनुष्यों में कोई
निएंडरथल मातृ रेखा नहीं बची है, चाहे इसलिए कि निएंडरथल माँओं के वंशज मिट गए
या इसलिए कि वह अंतःप्रजनन ज़्यादातर एक ही ओर चला। केंद्रकीय निशान बच गए; वह
वंशरेखा नहीं बची।

\textbf{16.} $1.2/0.27 \approx 4.4$ मिलियन साल --- यानी एक कम आकलन, क्योंकि
वह घड़ी अधिक दूर की गाँठों पर बिठाई गई थी।

\textbf{17.} वे जीवाश्म उस घड़ी के आकलन से पुराने हैं; असली गाँठ उस वंशरेखा
के सबसे पुराने जीवाश्म से कम से कम उतनी ही पुरानी है, यानी 6 से 7 मिलियन
साल, इसलिए इस हाल के अंतराल के लिए उस घड़ी की दर 0.27\% से कुछ धीमी होनी
चाहिए।

\textbf{18.} \emph{H. erectus} और, कालरेखा में, आख़िरी
\emph{Paranthropus} (तथा लगभग 0.7 से \emph{H. heidelbergensis});
इनमें से केवल उस वंशरेखा के जीवित वंशज हैं जो \emph{H. heidelbergensis} से
होकर \emph{H. sapiens} तक जाती है, और \emph{H. erectus} के केवल उन
\omterm{def:g12:selection-drift-speciation:population}{आबादियों} के रास्ते जो उसमें मिलीं --- यानी कई में से एक ही रेखा।

\textbf{19.} उस समय में कंकाल और मस्तिष्क शायद ही बदले, जबकि औज़ारों, कला,
भाषा, खेती और शहरों ने मनुष्यों के जीने का ढंग बदल डाला; वे बदलाव सीखने से
हस्तांतरित होते हैं, पीढ़ियों के भीतर जमा होते हैं, और अब उस पर्यावरण को
गढ़ते हैं जिसमें यह \omterm{def:g10:biodiversity-scales:species}{जाति} विकसित होती है।

\textbf{20.} चलने वाला A लगभग \qty{1.4}{m} का था और मनुष्य जैसे पैरों पर
धीरे चलता था; पैर पहले आए, मस्तिष्क एक मिलियन साल बाद, ठुड्डी आख़िर में; और
वह विद्यार्थिनी दो लुप्त \omterm{def:g12:selection-drift-speciation:population}{आबादियों} के \omterm{def:g10:universal-dna:gene}{जीन} ढोती है।
\end{solution}
